
\section{Vertex algebras, filtration, localisation and Lie algebras action} \label{section:vertex-algebra}

\subsection{Reminders on graded vertex algebras} \label{subsection:vertex-algebras}

For the details about the general theory of vertex algebras, see~\cite{kac1998vertex,frenkel2004vertex,desole2006finite}. We recall some basic notations and facts here.

Let $V$ be a vertex superalgebra over the complex number field $\C$. Denote by $1$ the vacuum vector, which is a non-zero even vector in $V$. Denote by $\partial : V \rightarrow V$ the translation operator, which is an even operator. The state-field correspondence is an even $\C$-linear homomorphism which maps $a \otimes b$ in $V \otimes V$ to a Laurent series $a(z) b$ in~$V(\!(z)\!)$. It is denoted by $a(z)b \coloneqq \sum_{n \in \Z} a_{(n)}b\, z^{-n-1}$ and the coefficient $a_{(n)}b$ is called the $n$-th product of $a$ and $b$. The formal power series
\[ a(z) \coloneqq \sum_{n \in \Z} a_{(n)} \quad \text{in }\quad (\operatorname{End} V)[\![z, z^{-1}]\!] \]
is called the \emph{field} associated to $a$.

The fundamental relation satisfied by the these fields is the Borcherds identity, see \cite[Section 3.3.10]{frenkel2004vertex} or \cite[(1.28)]{desole2006finite}:
\begin{equation} \label{equation:borcherds}
	\begin{multlined}
	\sum_{j \geqslant 0} (-1)^j \binom{n}{j} \big(a_{(m + n - j)}(b_{(k + j)} c) - (-1)^n (-1)^{|a||b|} b_{(n + k - j)} (a_{(m + j) c} \big) \\
	= \sum_{j \geqslant 0} \binom{m}{j} (a_{(n + j)} b)_{(m + k - j)} c,
	\end{multlined}  
\end{equation}
for any parity-homogeneous vectors $a, b, c$ in $V$ and any $m, n, k$ in $\Z$. We use the notation $|a|$ for the parity of $a$, that is to say $0$ if it is even, $1$ if it is odd.


We recall that the \emph{normally ordered product} is defined as $\NO{a b} \coloneqq a_{(-1)} b$. In particular, the vacuum satisfies $\NO{a 1} = \NO{a 1} = a$, and for any nonnegative integer $n$:
\[ \NO{\partial^n(a) b} = n! \, a_{(-n - 1)} b. \]
The normally ordered is non-associative in general, so we adopt the convention
\[ \NO{a b c} \coloneqq \NO{a \NO{b c}}, \]
where $a$, $b$ and $c$ are in $V$.

The \emph{operator product expansion (OPE)} of two elements $a$ and $b$ is the data of products $a_{(n)} b$ for all integers $n \geqslant 0$. It is denoted by:
\[ a(z_1) b(z_2) \sim \sum_{n \geqslant 0} \frac{1}{(z_1 - z_2)^n} (a_{(n)} b)(z_2). \]
Two element $a$ and $b$ are said to \emph{commute with each} other if and only if \[  [a(z_1), b(z_2)] = 0 \quad \text{in} \quad (\operatorname{End} V)[\![z_1, (z_1)^{-1}, z_2, (z_2)^{-1}]\!]. \]
By the Borcherds identity, it is equivalent to:
\[ a_{(n)} b = 0\quad \text{for all} \quad n \geqslant 0. \]


%
%For $a$ and $b$ in $V$, the {normally ordered product} is $\NO{a b} \coloneqq a_{(-1)} b$. For~$a, b, c$ in $\V$, we adopt the convention: $ \NO{a b c} \coloneqq  \NO{a (\NO{b c})}$. For $n$ a nonnegative integer, there is the relation $a_{(-n-1)} b = \frac{1}{n!} \NO{\partial(a) b}$. The $\lambda$-bracket is the formal sum defined by~$[a_\lambda b] \coloneqq \sum_{n=0}^\infty (a_{(n)}b)\frac{\lambda^n}{n!}$ in $V \otimes_\C \C[\lambda]$, it controls the nonassociativity and the noncommutativity of the normally ordered product.

In this paper, we call a \emph{graded vertex superalgebra} a vertex superalgebra $V$ such that there is a super vector space decomposition
\[ V = \bigoplus_{\Delta \in \Z_{\geqslant 0}} V(\Delta) \]
satisfying the inclusions
\[ V(\Delta){}_{(n)} V(\Delta') \subseteq V(\Delta + \Delta' - n - 1) \]
for any $n$ in $\Z$, and $\Delta, \Delta'$ in $\Z_{\geqslant 0}$. In particular $1$ must belong to $V(0)$ because $1_{(-1)}1 = 1$.

\begin{remark}
	It makes sense to consider a vertex algebra grading which is not supported on $\Z_{\geqslant 0}$. In this paper, such graded vertex algebras will not appear.
\end{remark}

\begin{theorem}[{\cite[Theorem 4.14]{li2004vertex}}] \label{theorem:li-filtration}
	Let $V$ be a graded vertex superalgebra. The data of subspaces
	\[ V_p \coloneqq \operatorname{Span} \Bigg\{ a_1{}_{(-n_1 - 1)} \cdots a_r{}_{(-n_r - 1)} 1 \quad \Bigg| \quad \begin{cases}
		r \in \Z_{\geqslant 0}, \\
		n_1, \dots, n_r \in \Z_{\geqslant 0}, \\
		\Delta_1, \dots, \Delta_r \in \Z_{\geqslant 0}, \\
		\Delta_1 + \cdots + \Delta_r \leqslant p, \\
		a_\bullet \in V(\Delta_\bullet)
	\end{cases} \Bigg\} \]
	for $p$ in $\Z_{\geqslant 0}$ defines an increasing, lower bounded and exhaustive filtration. For any integers $n$, $p$ and $q$, we have the embeddings:
	\begin{equation} \label{equation:li-filtration-property} \begin{aligned}
		V_p{}_{(n)} V_q & \subseteq V_{p + q} \quad \text{for} \quad n < 0, \\
		V_p{}_{(n)} V_q & \subseteq V_{p + q - 1} \quad \text{for} \quad n \geqslant 0.
	\end{aligned} \end{equation}
	Therefore, the associated graded vector space $\gr V$ inherits a super commutative vertex superalgebra structure from $V$ and a \emph{vertex Poisson structure}.
\end{theorem}

As in \cite{arakawa2026faithful}, we call this filtration the \emph{increasing Li filtration}. 


\begin{remark}
	In \cite[Definition 2.7 and Theorem 2.12]{li2005abelianizing}, Li introduces an upper bounded decreasing filtration $V_{\mathrm{Li}}^\bullet$, called the \emph{decreasing Li filtration} in \cite{arakawa2026faithful}, such that $\gr_{\mathrm{Li}} V$ is a vertex Poisson superalgebra. In fact, there is an isomorphism
	\[ \gr_{\mathrm{Li}} V \cong \gr V \]
	of vertex Poisson superalgebras \cite[Proposition 2.6.1]{arakawa2012remark}.
\end{remark}

%\begin{proposition}[{\cite[Proposition 2.6.1]{arakawa2012remark}}] \label{proposition:compare-filtrations}
%	Let $V$ be a graded vertex superalgebra. Then, for any integer $p$, $V_p$ and $V_{\mathrm{Li}}^p$ are graded subspaces and there is the equality
%	\[ V_{\mathrm{Li}}^p (\Delta) = V_{\Delta - p} (\Delta). \]
%	This equalities induce 
%\end{proposition} 

Let $\{a_i\}_{i=1}^n$ be a set of odd or even vectors in $V$. Denote by $\leqslant$ the lexicographic order on $\{1, \ldots, n\} \times \Z_{\geqslant 0}$. The vertex superalgebra $V$ is said \emph{freely generated} by the \emph{strong generators} $\{a_i\}_{i=1}^n$ if the products
\begin{equation*}
	a_{i_1}{}_{(-n_1 - 1)} \cdots a_{i_r}{}_{(-n_r - 1)} 1 \quad \text{where} \quad 
	\begin{cases}
		& r \geqslant 0 \quad \text{is an integer}, \\
		& i_\bullet \in \{1, \ldots, n\} \quad \text{and} \quad n_\bullet \in \Z_{\geqslant 0}, \\
		& (i_j, n_j) \leqslant (i_{j+1}, n_{j+1}) \quad \text{if} \quad a_{i_j} \quad \text{is even}, \\
		& (i_j, n_j) < (i_{j+1}, n_{j+1}) \quad \text{if} \quad a_{i_j} \quad \text{is odd},
	\end{cases}
\end{equation*}
form a basis of $V$ as a vector space. 

\subsection{Arc schemes and commutative vertex algebras} 
\label{subsection:arc-schemes}

Standard references for the theory of jet and arc schemes are \cite{ein2009jet,chambert2018motivic}. In the context of vertex algebras, see \cite[Section 2.1]{arakawa2015associated}.

Let $R$ be a (super)commutative unital associative superalgebra that is of finite type. Denote by $\J_\infty R$ the universal differential superalgebra generated by $R$, it comes with a natural superalgebra embedding $R \hookrightarrow \J_\infty R$ and satisfies the universal property given by \cite[Remark 3.1]{ein2009jet}.

\begin{lemma}[{\cite[Section 1.4]{frenkel2004vertex}}]
	Like any supercommutative differential algebra, $\J_\infty R$ is a supercommutative vertex algebra.
\end{lemma} 

Let $X \coloneqq \Spec R$ be the associated affine super scheme. The affine super scheme  $\J_\infty X \coloneqq \Spec \J_\infty R$ represents the functor 
\[ \text{superalgebra} ~ A \longmapsto \{\text{superalgebra homomorphisms} ~  R \rightarrow  A[\![t]\!]\}, \]
where $A[\![t]\!]$ is the superalgebra of formal power series with coefficients in $A$. The scheme $\J_\infty X$ is called the \emph{arc scheme} of $X$.

\begin{proposition}[{\cite[Proposition 2.3]{arakawa2012remark}}]
	If $R$ is a Poisson superalgebra, the Poisson bracket extends to a unique vertex Poisson super bracket on $\J_\infty R$.
\end{proposition}

For any nonnegative integer $m$, the functor
\[ \text{superalgebra} ~ A \longmapsto \{\text{superalgebra homomorphisms} ~  R \rightarrow  A[t]/(t^{m+1})\} \]
is represented by an affine scheme $\J_m X = \Spec \J_m R$, where $A[t]/(t^{m+1})$ denotes the algebra $A \otimes \C[t] / t^{m + 1} \C[t]$. The scheme $\J_m X$ is called the \emph{$m$-jet scheme} of $X$. They come with natural super scheme homomorphisms
\[ \J_\infty X \longrightarrow \J_{m+1} X \longrightarrow \J_{m} X, \]
called the truncation morphisms, making $\J_\infty X$ the projective limit of the $\J_m X$, for $m$ in $\Z_{\geqslant 0}$.

The jet and arc constructions commute with localisation in the following sense. If $S$ is a multiplicative subset of $R = \C[X]$, then there are natural algebra isomorphisms:
\[ \J_m(R[S^{-1}]) \cong (\J_mR)[S^{-1}] \quad \text{and} \quad \J_\infty(R[S^{-1}]) \cong (\J_\infty R)[S^{-1}].  \]

By construction, if $X$ is of finite type then $\J_m X$ is always of finite type for any finite $m$. But it is not the case for $\J_\infty X$ in general. If $X$ is a smooth irreducible variety then $\J_m X$ is a smooth irreducible variety for any finite $m$ \cite[Corollary~2.11]{chambert2018motivic}. 

If $N$ is an affine group scheme acting on an affine scheme $X$, then for any finite or infinite $m$, $\J_m N$ is an affine group scheme acting on $\J_m X$ \cite[Remark 2.8]{ein2009jet}. 

\subsection{Lie algebra cohomology for a projective limit of Lie algebras}

In this section, we review the properties of modules of the projective limit of a system of algebraic groups. This is closely related to the notion of smooth module introduced in \cite[Section 1.9]{kazhdan1993tensor}.

Let $\{N_m\}_{m \in \Z_{\geqslant 0}}$ be a projective system of algebraic groups such that each homomorphism $N_{m + 1} \rightarrow N_m$ is \emph{dominant}, meaning that $\C[N_m] \rightarrow \C[N_{m+1}]$ is \emph{injective}. Denote by $N_\infty$ the affine group scheme given by the projective limit.

\begin{example}
	If $N$ is an algebraic group then for any finite $m$, the jet scheme $\J_m N$ is an algebraic group. The arc group $\J_\infty N$ is a group scheme \cite[Remark 2.8]{ein2009jet} and is the projective limit of the algebraic groups $\J_m N$. If $N$ is connected then $\J_m N$ is connected. 
	
	The Lie algebras of the jet and arc groups are given by
	\[ \operatorname{\mathfrak{Lie}} \J_m N = \J_m \n \coloneqq \n[t] / (t^{m+1}) \quad \text{and} \quad \operatorname{\mathfrak{Lie}} \J_\infty N = \J_\infty \n \coloneqq \n[\![t]\!], \]
	where the brackets are induced by extension of scalars (it is a well-known fact, see \cite[Lemma 5.1.6.1]{juillard2025phd} for a proof).
\end{example}

\begin{proposition} \label{proposition:colimit-module-pro-group}
	Let $V$ be an $N_\infty$-module, and denote by $\rho \colon V \rightarrow \C[N_\infty] \otimes V$ the corresponding co-action map. Then the subspaces
	\[ V_{(m)} \coloneqq \{v \in V \quad | \quad \rho(v) \in \C[N_m] \otimes V \}, \quad \text{where} \quad m \in \Z_{\geqslant 0}, \] 
	define an inductive system of submodules such that, for each $m$, the action of $N_\infty$ on $V_{(m)}$ factorises through the natural morphism $N_\infty \rightarrow N_m$. Moreover, $V$ is the colimit of the $V_{(m)}$.
	
	If $\phi \colon V \rightarrow W$ is an $N_\infty$-module homomorphism, then for all integers $m$, there is an inclusion
	\[ \phi(V_{(m)}) \subseteq W_{(m)}. \]
\end{proposition}

\begin{remark}
	In the previous statement, when we write $\rho(v) \in \C[N_m] \otimes V$, we identify $\C[N_m]$ with its image in $\C[N_\infty]$.
\end{remark}

\begin{proof}
	We start with the first statement. If $v$ belongs to $V_{(m)}$, we can write $\rho(v)$ as a finite sum
	\[ \rho(v) = \sum_i F_i \otimes v_i \]
	where the $F_i$ are functions in $\C[N_m]$ and the $v_i$ are in $V$. Using the compatibility of the comodule structure with the counit and comultiplications on $\C[N_m]$ and $\C[N_\infty]$ \cite[Formulae (24)]{milne2017algebraic}, we get that 
	\[ v = \sum_i v_i \quad \text{and} \quad \rho(v_i) \in \C[N_m] \otimes V. \]
	Therefore, we deduce that
	\[ \rho(V_{(m)}) \subseteq \C[N_m] \otimes V_{(m)}, \]
	so $V_{(m)}$ is an $N_\infty$-submodule such that the action of $N_\infty$ factorises through the natural homomorphism $N_\infty \rightarrow N_m$. 
	
	There are obvious inclusions $V_{(m)} \subseteq V_{(m+1)}$ by using the group homomorphisms $N_{m+1} \rightarrow N_m$. 
	Therefore we get a natural module homomorphism
	\[ \operatorname{colim}\limits_{m \geqslant 0} V_{(m)} \hooklongrightarrow V. \]
	Let us prove it is surjective. For any $v$ in $V$, we can write $\rho(v)$ as a finite sum
	\[ \rho(v) = \sum_i F_i \otimes v_i, \] 
	where the $F_i$ are functions in $\C[N_\infty]$ and the $v_i$ are in $V$. Because we have a filtered colimit, these finitely many $F_i$ must all lie in one $\C[N_m]$ for some $m$ that is big enough and then $v$ belongs to this $V_{(m)}$. This concludes the proof of the isomorphism
	\[ \operatorname{colim}\limits_{m \geqslant 0} V_{(m)} \cong V. \]
	
	Let us consider the second statement. We have the equation
	\[ (\id_{\C[N_\infty]} \otimes \phi) \circ \rho = \rho \circ \phi, \]
	by definition of a $\C[N_\infty]$-comodule homomorphism $\phi \colon V \rightarrow W$. If $v$ belongs to $V_{(m)}$ then 
	\[ \rho(\phi(v)) =  (\id_{\C[N_\infty]} \otimes \phi) \circ \rho(v) =  (\id_{\C[N_m]} \otimes \phi) \circ \rho(v) \]
	belongs to $\C[N_m] \otimes W$, so $\phi(v)$ is in $W_{(m)}$.
\end{proof}

Recall that the action by left multiplication of an affine group scheme $N$ on itself defines an $N$-module structure on the ring of functions $\C[N]$.

\begin{lemma} \label{lemma:filtration-function-progroup}
	For any nonnegative integer $m$, the image of the algebra embedding $\C[N_m] \hookrightarrow \C[N_\infty]$ equals $\C[N_\infty]_{(m)}$.
\end{lemma}

\begin{proof}
	Denote by $\mu$ the multiplication on the group scheme $N_\infty$. The comultiplication
	\[ \mu^*\colon \C[N_\infty] \longrightarrow \C[N_\infty] \otimes \C[N_\infty] \]
	is the co-action in this case. Denote by $\varepsilon \colon \C[N_\infty] \rightarrow \C$ the algebra homomorphism defined by evaluating polynomial functions at the unit of the group. By \cite[Proposition 3.6]{milne2017algebraic}:
	\[ (\id_{\C[N_\infty]} \otimes \varepsilon) \circ \mu^*(F) = F \otimes 1 \]
	for all $F$ in $\C[N_\infty]$. If we take $F$ in $\C[N_\infty]_{(m)}$, we deduce that $F$ belongs to $\C[N_m]$ by applying the previous identity.
	
	This gives the inclusion $\C[N_\infty]_{(m)} \subseteq \C[N_m]$. The other inclusion holds because $N_\infty \rightarrow N_m$ is a group homomorphism, so we have an equality.
\end{proof}

Denote by $\n_m$ the Lie algebra of $N_m$ and $\n_\infty$ the Lie algebra of $N_\infty$, which is the projective limit of the $\n_m$. For $V$ an $N_\infty$-module, \Cref{proposition:colimit-module-pro-group} gives submodules $V_{(m)}$ with an induced action of $\n_m$. 

As explained in \Cref{subsection:filtered-representation}, we define the Lie algebra cohomology $\Hgy^\bullet(\n_m, V_{(m)})$ as the cohomology of the Chevalley--Eilenberg complex. The cochain complex given by the colimit
\[ \Cpx^\bullet(\n_\infty, V) \coloneqq \operatorname{colim}\limits_{m \geqslant 0} \Cpx^\bullet(\n_m, V_{(m)}) = V \otimes \operatorname{Alt}^\bullet\big(\operatorname{colim}\limits_{m \geqslant 0} (\n_m)^\vee\big) \]
is the Chevalley--Eilenberg complex for $V$. The Lie algebra cohomology of $V$ is then defined as
\[ \Hgy^\bullet(\n_\infty, V) \coloneqq \Hgy(\Cpx^\bullet(\n_\infty, V)) = \operatorname{colim}\limits_{m \geqslant 0}\Hgy^\bullet(\n_m, V_{(m)}), \]
where the last equality comes from the fact that cohomology commutes with filtered colimits.

\begin{lemma} \label{lemma:coholology-pro-algebra}
	Let $V$ be an $N_\infty$-module. Its degree $0$ Lie algebra cohomology gives the Lie algebra invariants:
	\[ \Hgy^0(\n_\infty, V) = V^{\n_\infty}. \]
	
	
	If all the $N_m$ are connected, then 
	\[ V^{N_\infty} = V^{\n_\infty}. \]
\end{lemma}

\begin{proof}
	The first statement is given by the following equalities:
	\[ V^{\n_\infty} =  \operatorname{colim}\limits_{m \geqslant 0} (V_{(m)})^{\n_m} = \operatorname{colim}\limits_{m \geqslant 0} \Hgy^0(\n_m, V_{(m)}) = \Hgy^0(\n_\infty, V). \]
	
	Because each $N_m$ is connected and an algebraic group (in particular a connected smooth variety), one has the equality $(V_{(m)})^{N_m} = (V_{(m)})^{\n_m}$. Taking the colimits, the second part of the statement follows.
\end{proof}

%Let $V$ be a $\n_\infty$-module. By analogy with the terminology of \cite{kazhdan1993tensor}, we say that $V$ is a smooth module if $V$ can be written as the colimit of an inductive system of submodules $V_{(m)}$, where $m$ is a nonnegative integer, such that the action of $\n_\infty$ on $V$ factorises through through the natural morphism $\n_\infty \rightarrow N_m$.

\begin{proposition} \label{proposition:gr-invariant-pro}
	Assume that, for all nonnegative integers $m$, $N_m$ is a graded unipotent group, such that its Lie algebra is nonpositively graded $\n_m = \bigoplus_{i \leqslant 0} \n_{m,i}$. Moreover, assume that the homomorphism $N_{m+1} \rightarrow N_m$ is compatible with the gradings. Then, the projective limit $N_\infty$ inherits a grading.
	
	Let $V$ be an $N_\infty$-module, equipped with an increasing filtration, that is filtered in the sense of \Cref{subsection:filtered-representation}. Make the following assumptions:
	\begin{enumerate}
		\item the filtration on $V$ is bounded from below and exhaustive,

		
		\item there is an isomorphism $\C[N_\infty] \otimes W \cong \gr V$ of $N_\infty$-modules, where $\C[N_\infty]$ is equipped with the left-multiplication action and $W$ denotes a trivial $N_\infty$-module.
	\end{enumerate}
	Then there is a natural graded linear isomorphism 
	$$ \gr (V^{N_\infty}) \overset{\sim}{\longrightarrow} (\gr V)^{N_\infty}. $$
	Moreover, the Lie algebra cohomology $\Hgy^n(\n_\infty, V)$ is zero whenever $n \neq 0$. 
\end{proposition}

\begin{proof}
	The argument is similar to the proof of \Cref{proposition:gr-invariant}. The cohomology of $V$ is defined as the cohomology of the complex
	\[ \Cpx^\bullet(\n_\infty, V) = V \otimes \operatorname{Alt}^\bullet\big(\operatorname{colim}\limits_{m \geqslant 0} (\n_m)^\vee\big). \]
	The grading on the Lie algebra and the filtration on $V$ equip this cochain complex with an exhaustive filtration which is bounded from below. Therefore, there is a spectral sequence~$\{E_{r}\}_{r=0}^\infty$ associated with this filtered cochain complex \cite[Section~5.4]{weibel1994introduction}.
	
	The $1$-page of this complex is given by the Lie algebra cohomology of $\gr V$:
	\[ \bigoplus_{p \in \Z} E_1^{p, p + n} = \Hgy^{n}(\n_\infty, \gr V) = \operatorname{colim}\limits_{m \geqslant 0} \Hgy^{n}(\n_m, (\gr V)_{(m)}), \]
	where $n$ is a nonnegative integer.
	
	By using \Cref{proposition:colimit-module-pro-group} and \Cref{lemma:filtration-function-progroup}, for any $m$ we get an $N_m$-equivariant isomorphism 
	\[ \C[N_m] \otimes W \cong (\gr V)_{(m)}. \]
	Therefore, we can apply \Cref{proposition:gr-invariant} and we get that
	\[ \Hgy^{n}(\n_m, (\gr V)_{(m)}) = 0 \quad \text{for} \quad n \neq 0. \]
	By taking the colimit, we get
	\[ \Hgy^{n}(\n_\infty, (\gr V)) = 0 \quad \text{for} \quad n \neq 0, \]
	so the spectral sequence $\{E_{r}\}_{r=0}^\infty$ collapses at the $1$-page. Therefore we can conclude the argument in the same way as for \Cref{proposition:gr-invariant}.
\end{proof}

\begin{example} \label{example:grading-jet-groups}
	If $N$ is a graded group, then the jet groups $\{\J_m N\}_{m \geqslant 0}$ form a projective system of graded groups. If $\n = \bigoplus_{i \in \Z} \n_i$ denotes the induced Lie algebra grading, the Lie algebra grading on $\n[t]/(t^{m+1})$ is the grading induced by $\n$ by imposing that the formal variable $t$ is of degree $0$.
\end{example}

\subsection{Localisation \emph{à la} Borisov} \label{subsection:localisation-vertex}

We explain the process of localisation for a graded vertex algebra as it was originally developed in \cite{borisov2001vertex,malikov1999chiral}. Our main reference is \cite[Sections 2 and 3]{arakawa2026faithful}.

In this subsection, let $V$ be a $\Z_{\geqslant 0}$-graded vertex algebra. The normally ordered product equips $V(0)$ with a commutative associative algebra structure with unit given by the vacuum vector. The embedding $V(0) \hookrightarrow V$ extends uniquely to a vertex algebra homomorphism
\[ \J_\infty V(0) \longrightarrow V, \]
making $V$ a vertex algebra module over $\J_\infty V(0)$ (that is called a $V(0)$-loop module in \cite{borisov2001vertex}).

\begin{theorem}[{\cite[Propositions 2.3 and 3.5]{borisov2001vertex} and \cite[Theorem 3.12]{arakawa2026faithful}}] \label{theorem:borisov-localisation}
	For any multiplicative subset $S$ in $V(0)$, one can define a graded vertex algebra $V[S^{-1}]$ equipped with a graded vertex algebra homomorphism $\iota \colon V \rightarrow V[S^{-1}]$ satisfying the following universal property. 
	
	For any graded vertex algebra $W$ and any graded vertex algebra homomorphism $\psi \colon V \rightarrow W$ such that $\psi(S)$ lies in the group of units $W(0)^\times$ of $W(0)$, there is a unique graded vertex algebra homomorphism $\widetilde{\psi} \colon V[S^{-1}] \rightarrow W$ such that $\widetilde{\psi} \circ \iota = \psi$.	
\end{theorem}

The vertex algebra $V[S^{-1}]$ is called the localisation of $V$ with respect to the multiplicative subset $S$. Such a localisation can also be defined for any vertex algebra module $M$ over $V$ that is (non-negatively) graded. One denotes by $\iota_M \colon M \rightarrow M[S^{-1}]$ this localisation. It satisfies an analogous universal property as $V[S^{-1}]$, see \cite[Proposition 2.3]{borisov2001vertex} and \cite[Proposition 3.2]{arakawa2026faithful}. Denote by $V\text{\rm{-}}\mathcal{mod}$ the category of graded $V$-modules.

\begin{proposition}[{\cite[Lemma 3.3]{arakawa2026faithful}}]
	Localisation defines a functor
	\[ V\text{\rm{-}}\mathcal{mod} \longrightarrow V[S^{-1}]\text{\rm{-}}\mathcal{mod}, \quad M \longmapsto M[S^{-1}] \]
	which is an exact functor.
\end{proposition}

Because $V[S^{-1}]$ is graded, it is equipped with an increasing Li filtration and one can show that this filtration is the localisation of the Li filtration on $V$ \cite[Equation (3.10)]{arakawa2026faithful}.

\begin{proposition}[{\cite[Lemma 3.13]{arakawa2026faithful}}]  \label{proposition:gr-borisov-localisation}
	There is a natural vertex Poisson algebra isomorphism
	\[ (\gr V)[S^{-1}] \overset{\sim}{\longrightarrow} \gr(V[S^{-1}]), \]
	where the left-hand side denotes the usual localised commutative algebra.
\end{proposition}

\subsection{Universal affine vertex algebra and its localisation}

To the simple Lie algebra $\g$ and a complex number $k$, one can associate a vertex algebra $\V^k(\g)$, called the \emph{universal affine vertex algebra}. See \cite[Section 2.1]{frenkel1992vertex}, \cite[Section 2.4]{frenkel2004vertex} or \cite[Section 1.8]{desole2006finite}. The vertex algebra $\V^k(\g)$ contains $\g$ as a linear subspace. It is freely generated by any ordered basis in $\g$ \cite[Proposition 1.20]{desole2006finite}.

If $\g = \bigoplus_{i \in \Z} \g_i$ is graded, then we can equip $\V^k(\g)$ with a vertex algebra grading
\[ \V^k(\g) = \bigoplus_{\Delta \in \Z} \V^k(\g)(\Delta) \] such that $\g_i$ is included in $\V^k(\g)(1 - i)$ \cite[Remark 1.23]{desole2006finite}. In general, such a grading is not necessarily supported on $\Z_{\geqslant 0}$, but if we pick our choice of good grading $\g = \g_{-1} \oplus \g_0 \oplus \g_1$ for the minimal nilpotent $f$, then we get that
\begin{equation} \label{equation:grading-universal-affine} 
	\V^k(\g) = \bigoplus_{\Delta \in \Z_{\geqslant 0}} \V^k(\g)(\Delta), 
\end{equation}
and the degree $0$ commutative subalgebra is:
\[ \Sym \g_1 \cong \V^k(\g)(0). \]
This can be seen by picking an ordered basis of $\g$ that is homogeneous for the grading.

By \Cref{theorem:li-filtration}, $\V^k(\g)$ is equipped with the increasing Li filtration and $\gr \V^k(\g)$ is a commutative vertex algebra. By using the fact that $\V^k(\g)$ is freely generated by any basis of $\g$, one can show that the inclusion of $\g$ in $\V^k(\g)$ induces a vertex Poisson algebra isomorphism
\begin{equation} \label{equation:gr-universal-affine}
	\C[\J_\infty \g^\vee] = \J_\infty(\Sym \g) \cong \gr \V^k(\g).
\end{equation}

Because $e$ belongs to $\g_1$, we can use the Borisov localisation process given by \Cref{theorem:borisov-localisation}, giving a graded vertex algebra homomorphism
\[ \V^k(\g) \longrightarrow \V^k(\g)[e^{-1}]. \]
We deduce from \eqref{equation:gr-universal-affine} and \Cref{proposition:gr-borisov-localisation} the vertex Poisson algebra isomorphism
\[ \J_\infty(\Sym \g)[e^{-1}] = \J_\infty(\Sym \g[e^{-1}]) \cong \gr (\V^k(\g)[e^{-1}]). \]

\begin{remark}
	The localisation homomorphism $\V^k(\g) \rightarrow \V^k(\g)[e^{-1}]$ is injective because its associated graded map $\J_\infty(\Sym \g) \rightarrow \J_\infty(\Sym \g)[e^{-1}]$ is injective, and the increasing Li filtration is bounded from below.
\end{remark}

Any module $M$ over $\V^k(\g)$ carries an action of the arc algebra $\J_\infty \g = \g[\![t]\!]$, that is given by the formula
\[ x t^n \cdot a = x_{(n)} a, \]
where $x$ is in $\g$, $n$ is a nonnegative integer and $a$ is in $M$. If one takes the vertex subalgebra $\V^k(\mathfrak h)$ generated by a Lie subalgebra $\mathfrak{h}$ in $\g$, then we get an action of $\J_\infty \mathfrak h = \mathfrak h[\![t]\!]$.

